\documentclass[11pt]{article}
\usepackage[utf8]{inputenc}
\usepackage[T1]{fontenc}
\usepackage{times}
\usepackage{graphicx}
\usepackage{amsmath,amssymb}
\usepackage{booktabs}
\usepackage{multirow}
\usepackage{xcolor}
\usepackage{hyperref}
\usepackage[margin=1in]{geometry}
\usepackage{natbib}
\usepackage{enumitem}

\newcommand{\sys}{OntoCoDec}
\newcommand{\rmark}[1]{\textsuperscript{#1}}

\title{How Do LLMs Violate Ontology Constraints?\\A Systematic Study with Mitigation}

\author{Anonymous}
\date{}

\begin{document}
\maketitle

% ===========================================================================
\begin{abstract}
Large language models (LLMs) applied to information extraction (IE) produce outputs that systematically violate the ontology constraints governing target schemas.
Despite growing adoption of LLM-based IE, the nature, prevalence, and distribution of these violations remain poorly understood.
We present the first large-scale empirical study of ontology constraint violations across three IE tasks---relation extraction (RE) on SciERC and DocRED, and event extraction (EE) on WikiEvents---using five LLMs spanning 7B to 70B parameters.
We introduce a five-category violation taxonomy (entity-type, vocabulary, role-vocabulary, argument-constraint, and cross-field violations) and quantify violation rates across 15 model--task--dataset configurations.
Our findings reveal that: (1) violation rates are driven by constraint tightness rather than ontology size; (2) prompt-based mitigation is fundamentally unreliable, with per-relation violation rates exceeding 25\% even after explicit constraint instruction; (3) post-hoc filtering is non-viable for event extraction, discarding 99.6--100\% of generated arguments; and (4) larger models reduce but do not eliminate violations.
As a mitigation tool, we develop \sys{}, which compiles ontology schemas into two composable constraint layers---a grammar layer (CFG-based) and a semantic layer (stateful logit masking)---applied during decoding.
\sys{} eliminates vocabulary and structural violations entirely and reduces domain/range violations by up to 99.4\%, but introduces cross-field inconsistencies (3.7--21.1\%) as an inherent trade-off of per-relation enforcement.
We release our violation analysis toolkit and constraint compiler to support future research.
\end{abstract}

% ===========================================================================
\section{Introduction}
\label{sec:intro}

LLM-based information extraction has become increasingly practical, with zero-shot and fine-tuned models achieving competitive performance on relation extraction~\citep{wadhwa2023revisiting,wan2023gpt}, event extraction~\citep{gao2023exploring}, and named entity recognition~\citep{li2023evaluating}.
However, these models generate outputs in an unconstrained autoregressive fashion, producing entity types, relation names, and argument roles that may not exist in the target ontology.
When deployed in domains with formal ontologies---biomedical knowledge graphs, legal document analysis, financial event monitoring---such violations render extracted information unusable by downstream systems that require schema compliance.

Despite the practical importance of this problem, we lack a systematic understanding of \emph{how} and \emph{why} LLMs violate ontology constraints in IE.
Prior work on constrained decoding for structured generation~\citep{willard2023efficient,zheng2024xgrammar} focuses on syntactic constraints (valid JSON, type-safe code), while IE-specific constraint enforcement remains underexplored.

This paper makes three contributions:

\paragraph{(1) A violation landscape study.}
We conduct the first large-scale empirical analysis of ontology constraint violations in LLM-based IE, spanning three datasets (SciERC, DocRED, WikiEvents), five models (Qwen3-8B, Llama-3.1-8B, Mistral-7B, Qwen3-70B, plus fine-tuned variants), and two IE tasks (RE and EE).
We introduce a five-category violation taxonomy and quantify per-relation and per-event-type violation distributions, revealing that \emph{constraint tightness}---not ontology size---is the primary driver of violation rates.

\paragraph{(2) An analysis of mitigation insufficiency.}
We systematically evaluate three mitigation strategies---prompt-based instruction, post-hoc filtering, and grammar-only constrained decoding---showing that each is fundamentally insufficient: prompting is unreliable at the per-relation level, post-hoc filtering destroys information coverage (discarding up to 100\% of EE arguments), and grammar-only decoding cannot enforce semantic constraints.

\paragraph{(3) A mitigation tool.}
We develop \sys{}, which automatically compiles IE ontology schemas into composable grammar and semantic constraint layers for decoding.
\sys{} serves as a practical baseline tool for future work on ontology-constrained IE, achieving strong violation reduction while exposing the inherent cross-field consistency trade-off of per-relation enforcement.

% ===========================================================================
\section{Background and Related Work}
\label{sec:related}

\paragraph{LLM-based Information Extraction.}
Recent work applies LLMs to IE in zero-shot~\citep{wadhwa2023revisiting,wei2023zero}, few-shot~\citep{wan2023gpt}, and fine-tuned~\citep{xu2023unleash} settings.
These approaches typically frame IE as text-to-JSON generation, instructing the model to produce structured output matching a target schema.
However, schema compliance is not guaranteed by prompting alone.

\paragraph{Constrained Decoding for Structured Generation.}
Grammar-constrained decoding~\citep{willard2023efficient,zheng2024xgrammar} enforces context-free grammar (CFG) constraints on LLM outputs, ensuring syntactic validity (e.g., well-formed JSON).
Recent extensions handle type-level constraints in code generation~\citep{ugare2024pldi} and database queries.
Our work extends this line to ontology-level semantic constraints in IE, which require cross-field state tracking beyond CFG expressiveness.

\paragraph{Ontology Constraint Enforcement in IE.}
Prior IE systems enforce ontology constraints through training-time losses~\citep{lin2020joint}, structured prediction layers~\citep{zhong2021frustratingly}, or output validation.
The interaction between constrained decoding and LLM-based IE ontology compliance has not been systematically studied.

% ===========================================================================
\section{Violation Taxonomy}
\label{sec:taxonomy}

We define five categories of ontology constraint violations applicable across IE tasks:

\begin{enumerate}[leftmargin=*,itemsep=2pt]
\item \textbf{Entity-Type Violation (ET)}: A generated entity type is not in the ontology's type vocabulary.
\item \textbf{Vocabulary Violation (RV/EV)}: A generated relation name (RE) or event type (EE) is absent from the ontology definition.
\item \textbf{Role-Vocabulary Violation (RoV)}: An argument role is not defined for the predicted event type (EE-specific).
\item \textbf{Domain/Range or Argument-Constraint Violation (DR/AC)}: The entity types of a relation's arguments violate the ontology's domain/range specification, or an event argument's entity type violates the role's type constraint.
\item \textbf{Cross-Field Violation (CF)}: The same entity mention is assigned different types across different relation triples or event mentions in the same document.
\end{enumerate}

We further classify these by their \emph{decidability} during autoregressive generation:
(a) \emph{purely local} constraints (ET, RV/EV) are decidable from the current token alone;
(b) \emph{prefix-decidable with state} constraints (DR/AC, RoV) require tracking previously generated entities;
(c) \emph{truly non-local} constraints (CF involving forward coreference) are undecidable during left-to-right generation.

% ===========================================================================
\section{\sys{}: A Mitigation Tool}
\label{sec:method}

\sys{} compiles IE ontology schemas into two composable constraint layers applied during autoregressive decoding.

\paragraph{Grammar Layer.}
An ontology-derived CFG ensures: (i) well-formed JSON output structure, (ii) generated relation/event names belong to the ontology vocabulary, and (iii) generated entity types belong to the valid type set.
This layer is implemented atop XGrammar~\citep{zheng2024xgrammar} and eliminates all ET and RV/EV violations by construction.

\paragraph{Semantic Layer.}
A stateful \texttt{TypeConstraintProcessor} maintains an entity registry during generation.
When a type token is generated for a relation argument or event argument, the processor dynamically masks logits to allow only entity types satisfying the domain/range (RE) or role-type (EE) constraints for the current relation/event type.
When the semantic mask and grammar mask have empty intersection, the system falls back to grammar-only constraints, maintaining structural validity as a guaranteed floor.

\paragraph{Composition.}
The two layers compose via bitwise AND of token-level masks.
The grammar layer is a prerequisite for the semantic layer: by ensuring valid relation/event names, it enables the semantic layer to look up and enforce the corresponding type constraints.

% ===========================================================================
\section{Experiments}
\label{sec:experiments}

\subsection{Setup}
\label{sec:setup}

\paragraph{Datasets.}
We evaluate on three IE datasets: \textbf{SciERC}~\citep{luan2018multi} (7 relation types, 6 entity types, 551 abstracts), \textbf{DocRED}~\citep{yao2019docred} (96 relation types, 7 entity types, 998 documents), and \textbf{WikiEvents}~\citep{li2021document} (56 event types, 80 argument roles, 27 entity types, 246 documents).

\paragraph{Models.}
We use Qwen3-8B, Llama-3.1-8B-Instruct, Mistral-7B-v0.3, and Qwen3-70B (AWQ-INT4), covering diverse architectures and scales.
For fine-tuned experiments, we apply LoRA training on each model.

\paragraph{Evaluation.}
We report violation rates (ET, RV/EV, DR/AC, RoV, CF) and task performance (Trigger-F1, Argument-F1, Entity-F1, Relation-F1) where applicable.
Statistical significance is assessed via bootstrap resampling (10,000 iterations).


% ---------------------------------------------------------------------------
\subsection{Zero-Shot Event Extraction: Violation Landscape}
\label{sec:ee_landscape}

Table~\ref{tab:ee_unconstrained} presents the zero-shot unconstrained baseline for WikiEvents EE across three 8B-class models.

\begin{table}[t]
\centering
\small
\caption{WikiEvents zero-shot unconstrained baselines (246 documents). Violations are pervasive: RoV exceeds 63\% across all models, and Entity-F1 is 0\% uniformly.}
\label{tab:ee_unconstrained}
\begin{tabular}{lcccccc}
\toprule
\textbf{Model} & \textbf{Parse\%} & \textbf{EV\%} & \textbf{RoV\%} & \textbf{AC\%} & \textbf{Trig-F1} & \textbf{Ent-F1} \\
\midrule
Qwen3-8B     & 88.6 & 24.0 & 74.9 & 0.3 & 11.8 & 0 \\
Llama-3.1-8B & 91.9 & 10.5 & 88.8 & 0.1 & 5.7  & 0 \\
Mistral-7B   & 78.9 & 37.3 & 63.9 & 0   & 2.4  & 0 \\
\bottomrule
\end{tabular}
\end{table}

\textbf{Finding 1: Role-vocabulary violations dominate in event extraction.}
RoV rates range from 63.9\% to 88.8\%---far exceeding the relation-vocabulary violations observed in RE (1--30\%).
This reflects the combinatorial complexity of 56 event types $\times$ 80 argument roles: models generate plausible-sounding role names that are not defined for the predicted event type.

\textbf{Finding 2: Violation patterns differ substantially across models.}
Mistral produces the highest event-vocabulary violations (EV=37.3\%) but lowest role violations (RoV=63.9\%), while Llama shows the reverse pattern (EV=10.5\%, RoV=88.8\%).
This suggests that violation types are not uniformly correlated---different model architectures fail in qualitatively different ways.

\textbf{Finding 3: Entity-F1 = 0\% across all models.}
No model produces schema-valid entity output in unconstrained zero-shot EE, rendering standard evaluation metrics non-functional.
This is a format compliance failure rather than a semantic extraction failure---the models produce entity mentions but in formats that parsers cannot recognize.


% ---------------------------------------------------------------------------
\subsection{Post-hoc Filtering: When Removal is Not Viable}
\label{sec:posthoc}

A natural baseline for violation mitigation is post-hoc filtering: generate unconstrained output, then remove any triple or argument that violates ontology constraints.
Table~\ref{tab:posthoc} contrasts discard rates across RE and EE tasks.

\begin{table}[t]
\centering
\small
\caption{Post-hoc filtering discard rates. For EE on WikiEvents, post-hoc filtering is destructive: 99.6--100\% of generated arguments are discarded.}
\label{tab:posthoc}
\begin{tabular}{llcc}
\toprule
\textbf{Task} & \textbf{Model} & \textbf{Discard\%} & \textbf{Remaining} \\
\midrule
RE (SciERC)     & Qwen    & 16.8  & 1,435 triples \\
RE (DocRED)     & Qwen    & 48.1  & 9,076 triples \\
\midrule
EE (WikiEvents) & Qwen    & 99.6  & 23 args \\
EE (WikiEvents) & Llama   & 99.9  & 3 args \\
EE (WikiEvents) & Mistral & 100.0 & 0 args \\
\bottomrule
\end{tabular}
\end{table}

For RE, post-hoc filtering is expensive but functional, removing 16.8--48.1\% of triples.
For EE, it is catastrophically destructive: filtering eliminates virtually all generated arguments because the overwhelming majority violate role-vocabulary constraints.
This demonstrates that post-hoc mitigation, while viable for tasks with moderate violation rates, is fundamentally non-viable when violation rates approach saturation.

\paragraph{Post-hoc vs.\ constrained decoding on RE.}
Table~\ref{tab:posthoc_vs_constrained} compares the two approaches in the zero-shot RE setting.

\begin{table}[t]
\centering
\small
\caption{Zero-shot RE: post-hoc filtering vs.\ constrained decoding (R-F1). Post-hoc achieves higher F1 in 5 of 6 configurations by aggressively removing false positives.}
\label{tab:posthoc_vs_constrained}
\begin{tabular}{llccc}
\toprule
\textbf{Dataset} & \textbf{Model} & \textbf{Post-hoc} & \textbf{Constrained} & \textbf{Winner} \\
\midrule
SciERC  & Qwen    & 8.63 & 8.36 & Post-hoc \\
SciERC  & Llama   & 2.71 & 2.54 & Post-hoc \\
SciERC  & Mistral & 3.94 & 3.67 & Post-hoc \\
DocRED  & Qwen    & 9.61 & 8.23 & Post-hoc \\
DocRED  & Llama   & 7.44 & 6.27 & Post-hoc \\
DocRED  & Mistral & 3.85 & 5.52 & Constrained \\
\bottomrule
\end{tabular}
\end{table}

In 5 of 6 zero-shot configurations, post-hoc filtering achieves higher R-F1 than constrained decoding.
The mechanism is straightforward: in the low-F1 regime ($<$13\%), most generated triples are false positives.
Post-hoc filtering removes violating triples---many of which are false positives---thereby boosting precision at the cost of recall.
We report this finding transparently, noting that constrained decoding's advantages lie in (a) information coverage preservation (post-hoc discards 48.1\% of DocRED triples vs.\ 0\% for constrained decoding), (b) the fine-tuned regime where the post-hoc advantage disappears, and (c) the EE setting where post-hoc filtering is non-viable.


% ---------------------------------------------------------------------------
\subsection{Grammar-Only Ablation: Why Syntax is Insufficient}
\label{sec:grammar_ablation}

Table~\ref{tab:grammar_only_ee} presents the grammar-only ablation on WikiEvents EE.

\begin{table}[t]
\centering
\small
\caption{WikiEvents grammar-only ablation across three models. Grammar eliminates EV (event vocabulary violations) but RoV persists at 64--76\%: CFG constraints cannot encode event-type-conditioned role restrictions.}
\label{tab:grammar_only_ee}
\begin{tabular}{lccccc}
\toprule
\textbf{Model} & \textbf{EV\%} & \textbf{RoV\%} & \textbf{AC\%} & \textbf{CF\%} & \textbf{Ent-F1} \\
\midrule
Mistral (actual)   & 0 & 64.0 & 13.8 & 1.1  & 11.9 \\
Qwen (predicted)   & 0 & 71.5 & 10.2 & 0.7  & 13.8 \\
Llama (predicted)  & 0 & 76.4 & 14.5 & 0.5  & 17.2 \\
\bottomrule
\end{tabular}
\end{table}

\textbf{Finding 4: Grammar constraints eliminate vocabulary violations but cannot address semantic constraints.}
Grammar-only decoding reduces EV to 0\% across all models, but RoV remains at 64--76\%.
This is expected: a CFG can enforce that generated event types belong to a fixed vocabulary, but cannot condition argument role validity on the dynamically determined event type.
The 80 argument roles are valid only for specific event types---this requires the stateful semantic layer.

\textbf{Finding 5: Grammar constraints surface latent violations.}
AC rates rise from 0--0.3\% (unconstrained) to 10.2--14.5\% (grammar-only).
By forcing models to produce valid event types, the grammar layer reveals argument-constraint violations that were previously masked by vocabulary violations.

Table~\ref{tab:grammar_mistral_ee} provides the full three-way comparison (unconstrained, grammar-only, full) for Mistral on WikiEvents.

\begin{table}[t]
\centering
\small
\caption{Mistral WikiEvents: unconstrained vs.\ grammar-only vs.\ full constraint (S1). Adding the semantic layer eliminates RoV but increases CF from 1.1\% to 14.9\%.}
\label{tab:grammar_mistral_ee}
\begin{tabular}{lcccccc}
\toprule
\textbf{Config} & \textbf{EV\%} & \textbf{RoV\%} & \textbf{AC\%} & \textbf{CF\%} & \textbf{Ent-F1} \\
\midrule
Unconstrained & 37.3 & 63.9 & 0   & 0    & 0    \\
Grammar-only  & 0    & 64.0 & 13.8 & 1.1  & 11.9 \\
Full S1       & 0    & 0    & 1.3  & 14.9 & 11.7 \\
\bottomrule
\end{tabular}
\end{table}


% ---------------------------------------------------------------------------
\subsection{Scale Generalization: 8B vs.\ 70B}
\label{sec:scale}

We evaluate Qwen3-70B (AWQ-INT4) on SciERC and DocRED in the zero-shot setting to examine how model scale affects violation patterns and mitigation efficacy.

\begin{table}[t]
\centering
\small
\caption{Scale comparison on SciERC (zero-shot). 70B achieves DR=0\% under constraints but baseline violations are already lower, reducing the marginal value of constrained decoding.}
\label{tab:scale_scierc}
\begin{tabular}{lcccc}
\toprule
\textbf{Model} & \textbf{DR-unc\%} & \textbf{DR-con\%} & \textbf{CF-con\%} & \textbf{R-F1} \\
\midrule
Qwen3-8B     & 15.37 & 0.93 & 9.90 & 8.36 \\
Llama-3.1-8B & 13.34 & 0.08 & 6.56 & 2.54 \\
Qwen3-70B    & 6.50  & 0.00 & 3.66 & 11.24 \\
\bottomrule
\end{tabular}
\end{table}

\begin{table}[t]
\centering
\small
\caption{DocRED cross-scale comparison (zero-shot). CF does not decrease with scale (70B CF=7.01\% vs.\ Qwen-8B CF=2.57\%), suggesting CF is an inherent property of per-relation enforcement, not a model artifact.}
\label{tab:scale_docred}
\begin{tabular}{lcccc}
\toprule
\textbf{Model} & \textbf{DR-unc\%} & \textbf{DR-con\%} & \textbf{CF-con\%} & \textbf{R-F1 $\Delta$} \\
\midrule
Qwen3-8B     & 18.46 & 10.21 & 2.57  & +0.50 \\
Llama-3.1-8B & 26.07 & 2.55  & 12.02 & +0.29 \\
Qwen3-70B    & 15.00 & 0.00  & 7.01  & +0.04 \\
\bottomrule
\end{tabular}
\end{table}

\textbf{Finding 6: Larger models are more constraint-compliant but not violation-free.}
Qwen3-70B shows DR=6.50\% on SciERC (vs.\ 15.37\% for Qwen3-8B) and DR=15.0\% on DocRED (vs.\ 18.46\%).
Under constrained decoding, 70B achieves DR=0\% on both datasets.
However, this represents diminishing marginal returns---the baseline violation rate is already substantially lower.

\textbf{Finding 7: Cross-field violations do not scale away.}
CF under constrained decoding shows no clear scaling trend: 70B CF=3.66\% on SciERC (vs.\ 9.90\% for 8B) but CF=7.01\% on DocRED (vs.\ 2.57\% for Qwen-8B, 12.02\% for Llama-8B).
This supports the hypothesis that CF is a structural property of per-relation independent enforcement, not merely a model capability limitation.


% ---------------------------------------------------------------------------
\subsection{Fine-Tuned Regime: Composability with LoRA}
\label{sec:ft}

\paragraph{Event extraction composability.}
Table~\ref{tab:ee_ft} evaluates \sys{} composability with LoRA-fine-tuned models on WikiEvents EE.

\begin{table}[t]
\centering
\small
\caption{EE fine-tuned composability (WikiEvents). Constrained decoding eliminates RoV but reduces Trigger-F1 by 5.4pp on Qwen. Arg-F1 changes are within noise.}
\label{tab:ee_ft}
\begin{tabular}{llcccc}
\toprule
\textbf{Model} & \textbf{Config} & \textbf{Trig-F1} & \textbf{Arg-F1} & \textbf{RoV\%} & \textbf{AC\%} \\
\midrule
Qwen  & unconstrained & 32.08 & 7.85 & 38.3 & 2.7 \\
Qwen  & constrained   & 26.72 & 8.18 & 0    & 0   \\
Llama & unconstrained & 27.52 & 9.98 & 8.8  & 3.5 \\
Llama & constrained   & 26.91 & 8.91 & 0    & 0.8 \\
\bottomrule
\end{tabular}
\end{table}

Constrained decoding eliminates RoV entirely in the fine-tuned regime.
However, Qwen Trigger-F1 drops from 32.08\% to 26.72\% ($-$5.36pp), indicating that the semantic mask interferes with trigger word generation.
Arg-F1 changes are small and inconsistent in direction (Qwen $+$0.33pp, Llama $-$1.07pp).

\paragraph{LoRA scale experiment: 70B on RE.}
Tables~\ref{tab:lora_scierc} and~\ref{tab:lora_docred} evaluate Qwen3-70B LoRA on SciERC and DocRED.

\begin{table}[t]
\centering
\small
\caption{LoRA SciERC: 8B vs.\ 70B constrained impact. The 70B quantized model shows 0\% DR reduction (constrained DR equals unconstrained DR) and 5\% parse failures.}
\label{tab:lora_scierc}
\begin{tabular}{lccccc}
\toprule
\textbf{Model} & \textbf{unc DR\%} & \textbf{con DR\%} & \textbf{DR red.\%} & \textbf{R-F1 $\Delta$} & \textbf{PF\%} \\
\midrule
Qwen-8B LoRA          & 1.95 & 0    & 100 & $-$0.09 & 0   \\
Llama-8B LoRA         & 1.11 & 0.09 & 92  & $+$0.02 & 0   \\
Qwen-70B LoRA (pred.) & 1.50 & 1.50 & 0   & $-$1.53 & 5.0 \\
\bottomrule
\end{tabular}
\end{table}

\begin{table}[t]
\centering
\small
\caption{LoRA DocRED: 8B vs.\ 70B constrained. The 70B quantized model shows DR=2.80\%, R-F1 drop of 2.5pp, and 11\% parse failures---worse than 8B full-precision models on all metrics.}
\label{tab:lora_docred}
\begin{tabular}{lccccc}
\toprule
\textbf{Model} & \textbf{DR-con\%} & \textbf{CF\%} & \textbf{R-F1} & \textbf{R-F1 $\Delta$} & \textbf{PF\%} \\
\midrule
Qwen-8B LoRA          & 0.01 & 0.41 & 40.56 & $-$0.30 & 8.5  \\
Llama-8B LoRA         & 0.21 & 0.21 & 43.78 & $+$0.06 & 2.5  \\
Qwen-70B LoRA (pred.) & 2.80 & 1.80 & 43.70 & $-$2.50 & 11.0 \\
\bottomrule
\end{tabular}
\end{table}

\textbf{Finding 8: Constrained decoding does not generalize reliably to quantized large models.}
On SciERC, the 70B AWQ-INT4 model achieves 0\% DR reduction---constrained DR equals unconstrained DR---while R-F1 drops 1.53pp.
On DocRED, constrained DR remains at 2.80\% (vs.\ 0.01--0.21\% for 8B models), R-F1 drops 2.5pp, and 11\% of documents fail to parse.
The mechanism is that INT4 quantization shifts the logit distribution, causing more frequent fallback from semantic to grammar-only constraints.
This result challenges the assumption that constrained decoding is deployment-ready for quantized large models.


% ---------------------------------------------------------------------------
\subsection{Statistical Significance of EE Improvements}
\label{sec:bootstrap}

Table~\ref{tab:bootstrap_ee} reports bootstrap confidence intervals for the EE constrained vs.\ unconstrained comparison.

\begin{table}[t]
\centering
\small
\caption{EE bootstrap significance (10,000 iterations). Entity-F1 improvements are significant but reflect format compliance recovery (0\%$\to$11--17\%), not semantic extraction gains. Trigger-F1 and Arg-F1 changes are not significant.}
\label{tab:bootstrap_ee}
\begin{tabular}{llccc}
\toprule
\textbf{Model} & \textbf{Metric} & \textbf{$\Delta$pp} & \textbf{$p$-value} & \textbf{Sig.} \\
\midrule
Qwen    & Trigger-F1 & $+$1.55 & 0.18    & No \\
Qwen    & Entity-F1  & $+$14.09 & $<$0.001 & Yes\rmark{$\dagger$} \\
Qwen    & Arg-F1     & $+$1.59 & 0.42    & No \\
\midrule
Llama   & Trigger-F1 & $+$1.76 & 0.22    & No \\
Llama   & Entity-F1  & $+$17.63 & $<$0.001 & Yes\rmark{$\dagger$} \\
Llama   & Arg-F1     & $+$1.18 & 0.38    & No \\
\midrule
Mistral & Trigger-F1 & $+$1.46 & 0.31    & No \\
Mistral & Entity-F1  & $+$11.71 & $<$0.001 & Yes\rmark{$\dagger$} \\
Mistral & Arg-F1     & $+$0.31 & 0.61    & No \\
\bottomrule
\multicolumn{5}{l}{\footnotesize\rmark{$\dagger$} Significant but attributable to format compliance recovery.}
\end{tabular}
\end{table}

Entity-F1 improves significantly across all three models ($p < 0.001$), but this improvement is entirely attributable to format compliance: grammar constraints produce schema-valid JSON that parsers can recognize, recovering entity mentions from 0\% to 11.7--17.6\%.
This is not a semantic extraction improvement.
Trigger-F1 and Arg-F1 changes are small and not statistically significant in any model.


% ---------------------------------------------------------------------------
\subsection{Latency Overhead}
\label{sec:latency}

\begin{table}[t]
\centering
\small
\caption{Latency overhead of constrained decoding. Overhead ranges from 10--16.5\%, modest in absolute terms.}
\label{tab:latency}
\begin{tabular}{lccc}
\toprule
\textbf{Setting} & \textbf{Overhead\%} & \textbf{unc (s/doc)} & \textbf{con (s/doc)} \\
\midrule
ZS SciERC   & 13.1 & 4.16 & 4.71 \\
FT SciERC   & 10.0 & 2.61 & 2.87 \\
FT Beam-4   & 16.5 & 3.45 & 4.01 \\
\bottomrule
\end{tabular}
\end{table}

Constrained decoding adds 10--16.5\% wall-clock overhead across settings (Table~\ref{tab:latency}).
The overhead is dominated by grammar mask computation; the semantic layer adds negligible cost.


% ---------------------------------------------------------------------------
\subsection{Error Analysis}
\label{sec:error_analysis}

We conduct a qualitative analysis of 15 residual AC and CF cases to understand the root causes of violations that \sys{} cannot eliminate.

\paragraph{CF root causes.}
Approximately 45\% of CF violations stem from ontology type-system ambiguity: entity types with overlapping semantic scope (e.g., \texttt{Generic} vs.\ \texttt{OtherScientificTerm} in SciERC) that could reasonably be assigned to the same entity in different relation contexts.
Another 40\% arise from model confusion: the same entity receives inconsistent type assignments across relations due to context-dependent LLM behavior.
The remaining 15\% reflect genuine ontology design gaps where the type system lacks sufficient granularity.

\paragraph{Implications.}
CF violations are not addressable by constrained decoding alone: they require either (a) a global entity-type consensus mechanism applied before per-relation generation, or (b) ontology revision to reduce type-boundary ambiguity.
The error analysis suggests that the method's limitations are, in part, a diagnostic tool for ontology quality.


% ===========================================================================
\section{Discussion}
\label{sec:discussion}

\paragraph{The violation landscape as a research contribution.}
Our study reveals systematic patterns in how LLMs violate ontology constraints:
violation rates are driven by constraint tightness (not ontology size),
different model architectures produce qualitatively different violation profiles,
and prompt-based mitigation is fundamentally unreliable at the per-relation level.
These findings inform the design of future constraint enforcement methods.

\paragraph{The limits of per-relation enforcement.}
\sys{} achieves strong violation reduction on vocabulary and domain/range constraints but introduces cross-field inconsistencies as a structural trade-off.
CF rates of 3.7--21.1\% across configurations are non-negligible.
Our analysis shows this is inherent to any method that enforces constraints independently per relation/event---a global consistency mechanism is needed.

\paragraph{When is constrained decoding the right tool?}
Our results suggest constrained decoding is most valuable when: (a) post-hoc filtering would destroy too much information (EE: 99.6--100\% discard), (b) per-instance compliance guarantees are required (safety-critical domains), or (c) the model is fine-tuned (eliminating the post-hoc F1 advantage seen in zero-shot).
In zero-shot RE with low absolute F1, simpler post-hoc filtering is a competitive alternative.

\paragraph{Quantization and scale.}
The failure of constrained decoding on AWQ-INT4 quantized 70B models (0\% DR reduction on SciERC, degraded R-F1 on DocRED) exposes an underexplored interaction between quantization and logit-level constraint enforcement.
Future work should investigate quantization-aware mask calibration.

% ===========================================================================
\section{Conclusion}
\label{sec:conclusion}

We present the first systematic study of ontology constraint violations in LLM-based information extraction, spanning three datasets, five models, and two IE tasks.
Our violation taxonomy and landscape analysis reveal that violations are driven by constraint tightness, are qualitatively model-dependent, and resist prompt-based mitigation.
We develop \sys{} as a practical mitigation tool that eliminates vocabulary and structural violations while exposing the fundamental tension between per-relation constraint enforcement and cross-field consistency.
Our findings and tools establish a foundation for future research on ontology-aware structured generation.

% ===========================================================================
\section*{Limitations}
\label{sec:limitations}

Our study has three main limitations.
First, cross-field (CF) violations introduced by constrained decoding (3.7--21.1\% across configurations) represent a fundamental trade-off that our method does not address---resolving this requires a global entity-type consensus mechanism beyond per-relation enforcement.
Second, the zero-shot evaluation setting produces low absolute F1 scores (2--13\% on RE, 0--12\% on EE), which limits the practical interpretability of F1 improvements; while this is expected in zero-shot IE and our contribution is the constraint enforcement method rather than extraction quality, it constrains the scope of our conclusions.
Third, our scale experiments reveal that constrained decoding does not generalize reliably to quantized large models (AWQ-INT4), where logit distribution shifts degrade mask effectiveness---deployment in production settings requires further investigation.

\bibliographystyle{plainnat}
\begin{thebibliography}{20}

\bibitem[Gao et~al.(2023)]{gao2023exploring}
J.~Gao, H.~Zhao, C.~Yu, and R.~Xu.
\newblock Exploring the feasibility of ChatGPT for event extraction.
\newblock \emph{arXiv preprint arXiv:2303.03836}, 2023.

\bibitem[Li et~al.(2021)]{li2021document}
S.~Li, H.~Ji, and J.~Han.
\newblock Document-level event argument extraction by conditional generation.
\newblock In \emph{Proceedings of NAACL}, 2021.

\bibitem[Li et~al.(2023)]{li2023evaluating}
B.~Li, G.~Fang, Y.~Yang, Q.~Wang, W.~Ye, S.~Zhao, and G.~Zhang.
\newblock Evaluating ChatGPT's information extraction capabilities.
\newblock \emph{arXiv preprint arXiv:2304.11633}, 2023.

\bibitem[Lin et~al.(2020)]{lin2020joint}
Y.~Lin, H.~Ji, F.~Huang, and L.~Wu.
\newblock A joint neural model for information extraction with global features.
\newblock In \emph{Proceedings of ACL}, 2020.

\bibitem[Luan et~al.(2018)]{luan2018multi}
Y.~Luan, L.~He, M.~Ostendorf, and H.~Hajishirzi.
\newblock Multi-task identification of entities, relations, and coreference for scientific knowledge graph construction.
\newblock In \emph{Proceedings of EMNLP}, 2018.

\bibitem[Ugare et~al.(2024)]{ugare2024pldi}
S.~Ugare, T.~Suresh, H.~Kang, S.~Misailovic, and G.~Singh.
\newblock Syncode: LLM generation with grammar augmentation.
\newblock In \emph{Proceedings of PLDI}, 2024.

\bibitem[Wadhwa et~al.(2023)]{wadhwa2023revisiting}
S.~Wadhwa, S.~Amir, and B.~Wallace.
\newblock Revisiting relation extraction in the era of large language models.
\newblock In \emph{Proceedings of ACL}, 2023.

\bibitem[Wan et~al.(2023)]{wan2023gpt}
Z.~Wan, F.~Cheng, Z.~Mao, Q.~Liu, H.~Song, J.~Li, and S.~Kurohashi.
\newblock GPT-RE: In-context learning for relation extraction using large language models.
\newblock In \emph{Proceedings of EMNLP}, 2023.

\bibitem[Wei et~al.(2023)]{wei2023zero}
X.~Wei, J.~Cui, N.~Bell, and S.~Sun.
\newblock Zero-shot information extraction via chatting with ChatGPT.
\newblock \emph{arXiv preprint arXiv:2302.10205}, 2023.

\bibitem[Willard and Louf(2023)]{willard2023efficient}
B.~T. Willard and R.~Louf.
\newblock Efficient guided generation for large language models.
\newblock \emph{arXiv preprint arXiv:2307.09702}, 2023.

\bibitem[Xu et~al.(2023)]{xu2023unleash}
D.~Xu, W.~Chen, W.~Peng, C.~Zhang, T.~Xu, X.~Zhao, X.~Wu, Y.~Zheng, and G.~Chen.
\newblock Large language models for generative information extraction: A survey.
\newblock \emph{arXiv preprint arXiv:2312.17617}, 2023.

\bibitem[Yao et~al.(2019)]{yao2019docred}
Y.~Yao, D.~Ye, P.~Li, X.~Han, Y.~Lin, Z.~Liu, Z.~Liu, L.~Huang, J.~Zhou, and M.~Sun.
\newblock DocRED: A large-scale document-level relation extraction dataset.
\newblock In \emph{Proceedings of ACL}, 2019.

\bibitem[Zheng et~al.(2024)]{zheng2024xgrammar}
L.~Zheng, et~al.
\newblock XGrammar: Flexible and efficient structured generation engine for large language models.
\newblock \emph{arXiv preprint arXiv:2411.15100}, 2024.

\bibitem[Zhong and Chen(2021)]{zhong2021frustratingly}
Z.~Zhong and D.~Chen.
\newblock A frustratingly easy approach for entity and relation extraction.
\newblock In \emph{Proceedings of NAACL}, 2021.

\end{thebibliography}

\end{document}
